\section{Enjambre de partículas}
\label{sec:pso}

% TODO: PSO ya está implementado en el notebook (run_pso, vectorizado en
% PyTorch), calibrado con Optuna y evaluado con el protocolo de 10 corridas.
% Falta redactar esta sección con la misma estructura que las Secciones 2 y 3:
%   - Formulación: v <- v + c1 r1 (p - x) + c2 r2 (g - x) y x <- x + v, sin
%     peso de inercia. El enjambre usa 5 partículas y el punto inicial
%     compartido es la posición inicial de la primera partícula.
%   - Calibración (punto 3a): Optuna ajusta c1 y c2 en [0.05, 2.5] con 25
%     iteraciones fijas. Reportar la tabla "Mejores hiperparámetros
%     encontrados con Optuna" y, como evidencia, las figuras de la
%     subsección "5. Curvas de aprendizaje de PSO con los mejores
%     hiperparámetros" del notebook.
%   - Protocolo de 10 corridas (punto 3b): mismos 10 puntos iniciales que
%     descenso del gradiente y RMSProp, máximo 50 iteraciones; usar las
%     tablas de la sección "Evaluación con los mejores hiperparámetros".
%     Ninguna corrida alcanza la tolerancia de 1e-3: sin inercia, las
%     partículas siguen oscilando alrededor del mejor global.
%   - Mejor corrida (punto 3c): puntos visitados por el enjambre sobre las
%     curvas de nivel, curva de aprendizaje y evolución de f por partícula.
%   - Respuesta a cómo se podría combinar descenso del gradiente con
%     simulated annealing y qué beneficios tendría (punto 3d).
%   - Pruebas unitarias documentadas en la Sección~\ref{sec:pruebas}.
